% Part: incompleteness
% Chapter: representability-in-q
% Section: minimization-representable

\documentclass[../../../include/open-logic-section]{subfiles}

\begin{document}

\olfileid{inc}{req}{min}
\olsection{नियमित लघुतमीकरण~$\Th{Q}$ मध्ये निरूपणीय असते}

अपरिबद्ध शोधाचा विचार करू. समजा $g(x, z)$ हे नियमित फलन
$\Th{Q}$ मध्ये, उदाहरणार्थ !!{formula}~$!A_g(x, z, y)$ या
सूत्राने, निरूपणीय आहे. $f(z) = \umin{x}{[g(x, z) = 0]}$
असे $f$ परिभाषित करू. $f$ चे निरूपण करणारे
!!a{formula}~$!A_f(z, y)$ शोधायचे आहे. $f(z)$ चे मूल्य
म्हणजे अशी संख्या~$x$ की (अ) $g(x, z) = 0$ आणि (आ) तशी
ती लघुतम आहे, म्हणजे कोणत्याही $w < x$ साठी
$g(w, z) \neq 0$. म्हणून पुढील निवड स्वाभाविक आहे:
\[
!A_f(z,y) \ident !A_g(y, z, \Obj 0) \land \lforall[w][(w < y \lif \lnot
  !A_g(w, z, \Obj 0))].
\]
अर्थात, सर्वसाधारण बाबतीत $z$ ऐवजी $z_0$, \dots, $z_k$
ही चलं घ्यावी लागतील.

या पुराव्यातही $\Th{Q}$ ची सिद्ध करण्याची शक्ती किती आहे,
याविषयी काही साहाय्यक विधाने लागतील.

\begin{lem}
\ollabel{lem:succ} प्रत्येक !!{constant}~$a$ आणि प्रत्येक नैसर्गिक
संख्या~$n$ साठी,
\[
\Th{Q} \Proves \eq[(a' + \num n)][(a + \num n)'].
\]
\end{lem}

\begin{proof}
नेहमीप्रमाणे $n$ वर विगमन वापरू. आधार टप्प्यात $n =
0$ असेल, तर $\Th{Q}$ हे $\eq[(a' + \Obj 0)][(a + \Obj
0)']$ सिद्ध करते, हे दाखवायचे आहे. पण पुढील विधाने मिळतात:
\begin{align}
  \Th{Q} & \Proves \eq[(a' + \Obj 0)][a'] \quad \text{स्वयंसिद्धक $Q_4$ ने}
  \ollabel{step1}\\
  \Th{Q} & \Proves  \eq[(a + \Obj 0)][a] \quad \text{स्वयंसिद्धक $Q_4$ ने}
  \ollabel{step2} \\
  \Th{Q} & \Proves \eq[(a + \Obj 0)'][a'] \quad \text{\olref{step2} वरून}
  \ollabel{step3} \\
  \Th{Q} & \Proves \eq[(a' + \Obj 0)][(a + \Obj 0)'] \quad
  \text{\olref{step1} आणि \olref{step3} वरून}\notag
\end{align}
विगमन टप्प्यात, $\Th{Q}
\Proves \eq[(a' + \num n)][(a + \num n)']$ हे सिद्ध झाले आहे
असे समजू. $\num{n+1}$ म्हणजे $\num{n}'$ असल्याने, आता
$\Th{Q}$ हे $\eq[(a' + \num{n}')][(a + \num{n}')']$
सिद्ध करते, हे दाखवायचे आहे. आपल्याला मिळते:
\begin{align}
  \Th{Q} & \Proves \eq[(a' + \num n')][(a' + \num n)'] \quad
  \text{स्वयंसिद्धक $!Q_5$ ने} \ollabel{step5}\\
  \Th{Q} & \Proves \eq[(a' + \num n)'][(a + \num n')'] \quad
  \text{विगमन गृहीतक, बेरीजेचे स्वयंसिद्धक आणि समतेचे नियम वापरून} \ollabel{step6}\\
  \Th{Q} & \Proves \eq[(a' + \num n')][(a + \num n')'] \quad
  \text{\olref{step5} आणि \olref{step6} वरून.} \notag
\end{align}
\end{proof}

हे पुन्हा लक्षात घ्यावे की यावरून $\Th{Q}$ हे
$\lforall[x][\lforall[y][(x' + y) = (x + y)']]$ सिद्ध करते,
असे म्हणता येत नाही; वरचे विधान त्याहून दुर्बल आहे. हे
!!{sentence}~$\Struct{N}$ मध्ये खरे असले, तरी $\Th{Q}$
त्याला सिद्ध करू शकत नाही.

\begin{lem}
\ollabel{lem:less-zero}
$\Th{Q} \Proves \lforall[x][\lnot x < \Obj 0]$.
\end{lem}

\begin{proof}
पुरावा अनौपचारिक रीतीने देऊ (म्हणजे औपचारिक !!{derivation}
कशी रचावी यासाठी फक्त संकेत देऊ).

कोणत्याही~$a$ साठी $\lnot a < \Obj 0$ हे सिद्ध करायचे आहे.
$<$ च्या व्याख्येनुसार $\Th{Q}$ मध्ये
$\lnot \lexists[y][\eq[(y'+a)][\Obj 0]]$ सिद्ध करायचे आहे.
$\lexists[y][\eq[(y'+a)][\Obj 0]]$ असे मानून व्याघात मिळवू.
$\eq[(b'+a)][\Obj 0]$ असे समजा. $!Q_3$ वापरल्यावर
$\eq[a][\Obj 0] \lor \lexists[y][\eq[a][y']]$ मिळते.
आता दोन बाबी पाहू.

पहिली बाब: $\eq[a][\Obj 0]$ खरे आहे. $\eq[(b'+a)][\Obj 0]$
वरून $\eq[(b' + \Obj 0)][\Obj 0]$ मिळते. $\Th{Q}$ च्या
स्वयंसिद्धक~$!Q_4$ ने $\eq[(b' + \Obj 0)][b']$ मिळते,
म्हणून $\eq[b'][\Obj 0]$ मिळते. पण स्वयंसिद्धक~$!Q_2$ ने
$\eq/[b'][\Obj 0]$ हेसुद्धा मिळते; हा व्याघात आहे.

दुसरी बाब: एखाद्या $c$ साठी $\eq[a][c']$. मग
$\eq[(b' +
c')][\Obj 0]$ मिळते. स्वयंसिद्धक~$!Q_5$ ने
$\eq[(b' + c)'][\Obj 0]$ मिळते; हेसुद्धा स्वयंसिद्धक~$Q_2$
शी विसंगत आहे.
\end{proof}

\begin{lem}
\ollabel{lem:less-nsucc}
प्रत्येक नैसर्गिक संख्या~$n$ साठी,
  \[
  \Th{Q} \Proves
  \lforall[x][(x < \num {n+1} \lif (\eq[x][\Obj 0] \lor \dots \lor
    \eq[x][\num n]))].
  \]
\end{lem}

\begin{proof}
$n$ वर विगमन वापरू. प्रथम $n = 0$ हा आधार टप्पा पाहू.
त्या वेळी कोणत्याही~$a$ साठी $a < \num 1 \lif \eq[a][\Obj 0]$
हे दाखवायचे आहे. $a < \num 1$ असे समजा. मग $<$ च्या
परिभाषक स्वयंसिद्धकाने $\lexists[y][\eq[(y'+a)][\Obj 0']]$
मिळते (कारण $\num 1 \ident
\Obj 0'$).

$b$ कडे हा गुणधर्म आहे, म्हणजे $\eq[(b'+a)][\Obj 0']$,
असे समजा. $\eq[a][\Obj 0]$ दाखवायचे आहे. स्वयंसिद्धक~$!Q_3$
ने $\eq[a][\Obj 0]$ किंवा एखाद्या $c$ साठी
$\eq[a][c']$ मिळते. पहिल्या बाबतीत आणखी काही दाखवायचे नाही.
म्हणून $\eq[a][c']$ असे समजा. मग
$\eq[(b' + c')][\Obj 0']$ मिळते. $\Th{Q}$ च्या स्वयंसिद्धक~$!Q_5$
ने $\eq[(b'+c)'][\Obj 0']$ मिळते. स्वयंसिद्धक~$!Q_1$ ने
$\eq[(b' + c)][\Obj
0]$ मिळते. पण स्वयंसिद्धक~$!Q_8$ नुसार याचा अर्थ
$c < \Obj 0$; याला \olref{lem:less-zero} ने विरोध होतो.

आता विगमन टप्पा पाहू. $n$ साठीचे विधान गृहीत धरून $n+1$
साठी ते सिद्ध करू. $a < \num {n+2}$ असे समजा. पुन्हा
$!Q_3$ वापरून दोन बाबी मिळतात: $\eq[a][\Obj 0]$, किंवा
एखाद्या $b$ साठी $\eq[a][b']$. पहिल्या बाबतीत
$\eq[a][\Obj 0] \lor \dots \lor
\eq[a][\num{n+1}]$ सहज मिळते. दुसऱ्या बाबतीत $b'
< \num {n+2}$, म्हणजे $b' < \num{n+1}'$. स्वयंसिद्धक~$!Q_8$
नुसार एखाद्या $c$ साठी
$\eq[(c'+b')][\num{n+1}']$. स्वयंसिद्धक $!Q_5$ ने
$\eq[(c'+b)'][\num{n+1}']$. स्वयंसिद्धक~$!Q_1$ ने
$\eq[(c'+b)][\num{n+1}]$; म्हणून स्वयंसिद्धक~$!Q_8$ ने
$b < \num{n+1}$. विगमन गृहीतकाने
$\eq[b][\Obj 0] \lor \dots \lor \eq[b][\num{n}]$.
यावरून तर्कशास्त्राने
$\eq[b'][\Obj 0'] \lor \dots \lor \eq[b'][\num{n}']$
मिळते; आणि $\eq[a][b']$ असल्याने
$\eq[a][\num{1}] \lor \dots \lor \eq[a][\num{n+1}]$
मिळते. उलट दिशेसाठी, उजवीकडील प्रत्येक निर्दिष्ट संख्यांकाशी
समता गृहीत धरल्यास, वरच्या मर्यादेपर्यंतचे उरलेले अंतर
दाखवणारी विशिष्ट संख्या साक्षी म्हणून घ्यावी. आधी सिद्ध
केलेल्या संख्यांकांच्या बेरीजविषयक विधानाने अपेक्षित बेरीज
सिद्ध होते आणि क्रमसंबंधाच्या परिभाषक स्वयंसिद्धकाने तो संख्यांक
वरच्या मर्यादेपेक्षा लहान ठरतो;
आधार टप्प्यातही याच रीतीने उलट दिशा मिळते.
\end{proof}

\begin{lem}
  \ollabel{lem:trichotomy} प्रत्येक नैसर्गिक संख्या~$m$ साठी,
  \[
  \Th{Q} \Proves
  \lforall[y][((y < \num{m} \lor \num{m} < y) \lor \eq[y][\num{m}])].
  \]
\end{lem}

\begin{proof}
$m$ वर विगमन वापरू. प्रथम $m=0$ ही बाब पाहू. $!Q_3$ ने
$\Th{Q} \Proves
\lforall[y][(\eq[y][\Obj 0] \lor \lexists[z][\eq[y][z']])]$.
$a$ कोणताही असू द्या. मग $\eq[a][\Obj 0]$ किंवा एखाद्या~$b$
साठी $\eq[a][b']$. पहिल्या बाबतीत
$(a < \Obj 0 \lor \Obj
0 < a) \lor \eq[a][\Obj 0]$ हेसुद्धा मिळते. पण
$\eq[a][b']$ असेल, तर $\eq$ च्या तर्कशास्त्राने
$\eq[(b' +
\Obj 0)][(a + \Obj 0)]$ मिळते. $!Q_4$ ने
$\eq[(a +
\Obj 0)][a]$, म्हणून $\eq[(b' + \Obj 0)][a]$ आणि त्यामुळे
$\lexists[z][\eq[(z' + \Obj 0)][a]]$ मिळते. $!Q_8$
मधील $<$ च्या व्याख्येनुसार $\Obj 0 < a$. $\Obj 0 < a$
असेल, तर $(\Obj 0 < a \lor a
< \Obj 0) \lor \eq[a][\Obj 0]$ हेसुद्धा मिळते.

आता पुढील विधान गृहीत धरू:
\begin{align*}
  \Th{Q} & \Proves \lforall[y][((y < \num{m} \lor \num{m} < y) \lor
    \eq[y][\num{m}])]
  \intertext{आणि पुढील विधान सिद्ध करायचे आहे:}
  \Th{Q} & \Proves \lforall[y][((y < \num{m+1} \lor \num{m+1} < y) \lor
    \eq[y][\num{m+1}])]
\end{align*}
$a$ कोणताही असू द्या. $!Q_3$ ने $\eq[a][\Obj 0]$ किंवा
एखाद्या~$b$ साठी $\eq[a][b']$. पहिल्या बाबतीत $!Q_4$ ने
$\eq[\num{m}' +
a][\num{m+1}]$ मिळते; म्हणून $!Q_8$ ने
$a < \num{m+1}$.

आता दुसरी बाब, $\eq[a][b']$, पाहू. विगमन गृहीतकानुसार
$(b < \num{m} \lor \num{m} < b) \lor
    \eq[b][\num{m}]$.

पहिला विकल्प $b < \num{m}$ हा $!Q_8$ नुसार
$\lexists[z][\eq[(z' + b)][\num{m}]]$ याच्याशी समतुल्य आहे.
एखाद्या $c$ कडे हा गुणधर्म आहे असे समजा.
$\eq[(c' + b)][\num{m}]$ असेल, तर
$\eq[(c' + b)'][\num{m}']$ हेसुद्धा मिळते. $!Q_5$ ने
$\eq[(c' + b)'][(c' + b')]$ मिळते. म्हणून
$\eq[(c'+b')][\num{m}']$ मिळते. $\num{m}'
\ident \num{m+1}$ हे लक्षात ठेवून $c$ वर अस्तित्ववाचक
सामान्यीकरण केल्यावर
$\lexists[u][\eq[(u' + b')][\num{m+1}]]$ मिळते.
म्हणून $b < \num{m}$ असेल, तर $b' < \num{m+1}$, आणि
त्यामुळे $a < \num{m+1}$.

आता $\num{m} < b$ असे समजा, म्हणजे
$\lexists[z][\eq[(z' +
\num{m})][b]]$. $c$ हा असा~$z$ असेल, म्हणजे
$\eq[(c' +
\num{m})][b]$, असे समजा. तर्कशास्त्राने
$\eq[(c' + \num{m})'][b']$ मिळते. $!Q_5$ ने
$\eq[(c' + \num{m}')][b']$ मिळते. $\eq[a][b']$ आणि
$\num{m}' \ident
\num{m+1}$ असल्याने $\eq[(c' + \num{m+1})][a]$.
$!Q_8$ ने $\num{m+1} < a$.

शेवटी $\eq[b][\num{m}]$ असे समजू. मग तर्कशास्त्राने
$\eq[b'][\num{m}']$, आणि म्हणून
$\eq[a][\num{m+1}]$ मिळते.

अशा प्रकारे $m$ आणि~$b$ यांच्या बाबतीतील प्रत्येक विकल्पापासून
$m+1$ आणि~$a$ यांच्या बाबतीतील अनुरूप विकल्प मिळतो.
\end{proof}

\begin{prop}
\ollabel{prop:rep-minimization}
$!A_g(x, z, y)$ हे $g(x, z)$ चे~$\Th{Q}$ मध्ये निरूपण
करत असेल आणि ते फलन नियमित असेल, तर
\[
!A_f(z,y) \ident !A_g(y, z, \Obj 0) \land \lforall[w][(w < y \lif \lnot
  !A_g(w, z, \Obj 0))]
\]
हे सूत्र $f(z) = \umin{x}{[g(x, z) = 0]}$ चे निरूपण करते.
\end{prop}

\begin{proof}
प्रथम $f(n) = m$ असेल, तर
$\Th{Q} \Proves !A_f(\num n, \num m)$, म्हणजे पुढील विधान,
हे दाखवू:
\begin{align}
  \Th{Q} & \Proves !A_g(\num{m}, \num{n}, \Obj 0) \land \lforall[w][(w <
    \num{m} \lif \lnot !A_g(w, \num{n}, \Obj 0))]. \notag
  \intertext{$!A_g(x, z, y)$ हे $g(x, z)$ चे निरूपण करते आणि
    $f(n) = m$ असेल तेव्हा $g(m, n) =
    0$; म्हणून}
\Th{Q} & \Proves !A_g(\num{m}, \num{n}, \Obj 0). \notag
\intertext{$f(n) = m$ असेल, तर प्रत्येक $k < m$ साठी
  $g(k, n) \neq 0$. निरूपणातील एकमेव-मूल्याची अट आणि
  भिन्न संख्यांकांची असमता वापरल्यावर}
\Th{Q} & \Proves \lnot !A_g(\num{k}, \num{n}, \Obj 0). \notag
\intertext{यावरून पुढील विधान मिळते:}
\Th{Q} & \Proves \lforall[w][(w < \num{m} \lif \lnot
  !A_g(w, \num{n}, \Obj 0))]. \ollabel{rep-less}
\end{align}
$m = 0$ असेल, तर हे \olref{lem:less-zero} वरून;
अन्यथा \olref{lem:less-nsucc} वरून मिळते.

आता $f(n) = m$ असेल, तर
$\Th{Q} \Proves
\lforall[y][(!A_f(\num{n}, y) \lif \eq[y][\num{m}])]$
हे दाखवू. औपचारिक सिद्धता वाचकावर सोडून युक्तिवाद पुन्हा
अनौपचारिकपणे सांगू.

$!A_f(\num{n}, b)$ असे समजा. त्यावरून (अ)
$!A_g(b, \num{n},
\Obj 0)$ आणि (आ) $\lforall[w][(w < b \lif \lnot !A_g(w, \num{n}, \Obj
  0))]$ मिळते. \olref{lem:trichotomy} नुसार
$(b < \num{m} \lor \num{m} < b)
\lor \eq[b][\num{m}]$. $b < \num{m}$ किंवा $\num{m}
< b$ या दोन्ही शक्यतांपासून व्याघात मिळतो, हे दाखवू.

$\num{m} < b$ असेल, तर (आ) वरून
$\lnot !A_g(\num{m}, \num{n}, \Obj 0)$ मिळते. पण
$m = f(n)$ असल्याने $g(m, n) = 0$; आणि $!A_g$ हे~$g$
चे निरूपण करत असल्याने $\Th{Q} \Proves
!A_g(\num{m}, \num{n}, \Obj 0)$. हा व्याघात आहे.

आता $b < \num{m}$ असे समजा. \olref{rep-less} नुसार
$\Th{Q} \Proves \lforall[w][(w <
  \num{m} \lif \lnot !A_g(w, \num{n}, \Obj 0))]$,
म्हणून $\lnot !A_g(b, \num{n}, \Obj 0)$ मिळते.
हे पुन्हा (अ) शी विसंगत आहे. उलट अभिव्यंजन आधी सिद्ध
केलेल्या निरूपक सूत्राच्या निर्दिष्ट मूल्याच्या उदाहरणावरून आणि
समतेच्या प्रतिस्थापन नियमावरून मिळते.
\end{proof}

\end{document}
